% 来源及取材范围见分题文件；此为整理者转录的正文，不是下载原件。
\begin{lemma}[11.4.6, Tag 074E]
Let $A$ be a finite simple $k$-algebra.
\begin{enumerate}
\item There exists exactly one simple $A$-module $M$ up to isomorphism.
\item Any finite $A$-module is a direct sum of copies of a simple module.
\item Two finite $A$-modules are isomorphic if and only if they have the same dimension over $k$.
\item If $A=\operatorname{Mat}(n\times n,K)$ with $K$ a finite skew field extension of $k$, then $M=K^{\oplus n}$ is a simple $A$-module and $\operatorname{End}_A(M)=K^{op}$.
\item If $M$ is a simple $A$-module, then $L=\operatorname{End}_A(M)$ is a skew field finite over $k$ acting on the left on $M$, we have $A=\operatorname{End}_L(M)$, and the centers of $A$ and $L$ agree. Also $[A:k][L:k]=\dim_k(M)^2$.
\item For a finite $A$-module $N$ the algebra $B=\operatorname{End}_A(N)$ is a matrix algebra over the skew field $L$ of (5). Moreover $\operatorname{End}_B(N)=A$.
\end{enumerate}
\end{lemma}
\begin{proof}
By Theorem 11.3.3 we can write $A=\operatorname{Mat}(n\times n,K)$ for some finite skew field extension $K$ of $k$. By Lemma 11.4.5 the category of modules over $A$ is equivalent to the category of modules over $K$. Thus (1), (2), and (3) hold because every module over $K$ is free. Part (4) holds because the equivalence transforms the $K$-module $K$ to $M=K^{\oplus n}$. Using $M=K^{\oplus n}$ in (5) we see that $L=K^{op}$.

The statement about the center of $L=K^{op}$ follows from Lemma 11.4.5. The statement about $\operatorname{End}_L(M)$ follows from the explicit form of $M$. The formula of dimensions is clear. Part (6) follows as $N$ is isomorphic to a direct sum of copies of a simple module.
\end{proof}

\begin{lemma}[11.4.7, Tag 074F]
Let $A$, $A'$ be two simple $k$-algebras one of which is finite and central over $k$. Then $A\otimes_k A'$ is simple.
\end{lemma}
\begin{proof}
Suppose that $A'$ is finite and central over $k$. Write $A'=\operatorname{Mat}(n\times n,K')$, see Theorem 11.3.3. Then the center of $K'$ is $k$ and we conclude that $A\otimes_k K'$ is simple by Lemma 11.4.4. Hence $A\otimes_k A'=\operatorname{Mat}(n\times n,A\otimes_k K')$ is simple by Lemma 11.4.5.
\end{proof}

\begin{theorem}[11.6.1, Skolem--Noether, Tag 074Q]
Let $A$ be a finite central simple $k$-algebra. Let $B$ be a simple $k$-algebra. Let $f,g:B\to A$ be two $k$-algebra homomorphisms. Then there exists an invertible element $x\in A$ such that $f(b)=xg(b)x^{-1}$ for all $b\in B$.
\end{theorem}
\begin{proof}
Choose a simple $A$-module $M$. Set $L=\operatorname{End}_A(M)$. Then $L$ is a skew field with center $k$ which acts on the left on $M$, see Lemmas 11.3.2 and 11.4.6. Then $M$ has two $B\otimes_k L^{op}$-module structures defined by $m\cdot_1(b\otimes l)=lmf(b)$ and $m\cdot_2(b\otimes l)=lmg(b)$. The $k$-algebra $B\otimes_k L^{op}$ is simple by Lemma 11.4.7. Since $B$ is simple, the existence of a $k$-algebra homomorphism $B\to A$ implies that $B$ is finite.

Thus $B\otimes_k L^{op}$ is finite simple and we conclude the two $B\otimes_k L^{op}$-module structures on $M$ are isomorphic by Lemma 11.4.6. Hence we find $\varphi:M\to M$ intertwining these operations. In particular $\varphi$ is in the commutant of $L$ which implies that $\varphi$ is multiplication by some $x\in A$, see Lemma 11.4.6. Working out the definitions we see that $x$ is a solution to our problem.
\end{proof}
